\documentclass[10pt,a4paper]{amsart}
\usepackage[margin=19mm]{geometry}
\usepackage{amsmath,amssymb,mathtools}
\usepackage[scr=rsfs,cal=euler]{mathalfa}
\usepackage{xfrac}
\usepackage[unicode,hidelinks,bookmarks=false]{hyperref}
\newcommand{\ie}{i.e.\ }
\newcommand{\cf}{cf.\ }
\newcommand{\resp}{respectively\ }
\newcommand{\Subsection}{\subsection}
\providecommand{\nobreakdash}{\nobreak\hbox{-}}
\setlength{\emergencystretch}{2em}
\multlinegap0pt
\usepackage{amssymb}
\usepackage{mathtools}
\usepackage{xcolor}
\usepackage[shortlabels]{enumitem}
\newtheorem{proposition}{Proposition}
\newtheorem{theorem}{Theorem}
\newcommand{\F}{\mathbb{F}}
\newcommand{\rs}{\mathrm{rowspace}\,}
\title{On set-like sunflower-free families of subspaces\\over finite fields}
\author{Kamil Otal}
\date{Source: arXiv:2605.12232v1 (CC BY 4.0)}
\begin{document}
\maketitle

\noindent\emph{Source and licence.} On set-like sunflower-free families of subspaces\\over finite fields. Version arXiv:2605.12232v1 (CC BY 4.0), distributed by arXiv under the Creative Commons Attribution 4.0 International licence, http://creativecommons.org/licenses/by/4.0/, as recorded in the arXiv licence metadata for this exact version and shown on the abstract page https://arxiv.org/abs/2605.12232v1. Full licence notice: \texttt{licenses/arxiv-2605.12232-CC-BY-4.0.txt}.\par\smallskip

\noindent\textit{Editorial scope.} Counted unit: the article's theorem thm:construction (source0.tex lines 204--210). The article's complete original proof of it is delivered with it (source0.tex lines 212--244, one \texttt{proof} environment of 33 lines). No mathematics is written, repaired or re-derived: the statement and the proof are the article's own lines, in the article's own notation, and no retained line is substituted or re-flowed.  Retained uncounted as the source-local context the proof needs: lines 120--122; lines 195--200.  Nothing was omitted from any retained span, so the package carries no omission record.  No retained line contains a citation, so the wrapper carries no bibliography and no bibliography entry.  No retained line contains a figure environment, an image-inclusion command or drawing code, so no image is delivered and the package declares no asset dependency.  Editorial formatting only, in addition: an AMS \texttt{amsart} preamble that restates the article's own declarations and macros used by the retained lines, namely the environments \texttt{proposition}, \texttt{theorem} on separate counters, together with the standard packages those lines need.  The retained environments print in this document as Proposition 1, Theorem 1 in order of appearance, whereas the article prints the same items with the numbering of its own full document.  The complete native source of the article as distributed in the arXiv e-print for this exact version is bundled unchanged as \texttt{sources/200411/source0.tex}.
\begin{proposition}\label{prop:matrix-rep}
Let $\mathcal{C} \subseteq \F_q^{k \times k}$ be the matrix representation of $\F_{q^k}$. Then $\mathcal{C}$ is an MRD code with $d_R(\mathcal{C}) = k$.
\end{proposition}

\begin{proposition}\label{prop:sum-complement}
Let $U_1, U_2, \ldots, U_s$ be subspaces of $\F_q^n$. Then
\[
(U_1 + U_2 + \cdots + U_s)^\perp \;=\; U_1^\perp \cap U_2^\perp \cap \cdots \cap U_s^\perp.
\]
\end{proposition}

\begin{theorem}\label{thm:construction}
Let $n \geq 2\ell + 1$ and $k = n - \ell$. Let $I$ denote the $\ell \times \ell$ identity matrix over $\F_q$, and let $\mathcal{C} \subseteq \F_q^{\ell \times \ell}$ be the matrix representation of $\F_{q^\ell}$ over $\F_q$, as in Proposition~\ref{prop:matrix-rep}. For a matrix $A$, write $[A]_1$ for its first column, and let $\mathbf{0}$ denote the $\ell \times (n - 2\ell - 1)$ zero matrix. Define
\[
\mathcal{G} \;=\; \bigl\{ \rs [I \mid A \mid [A^2]_1 \mid \mathbf{0}] : A \in \mathcal{C}\bigr\}.
\]
Then $\mathcal{F} = \{U^\perp : U \in \mathcal{G}\}$ is a set-like $s$-sunflower-free family of $k$-spaces of $\F_q^n$ of size $q^\ell = q^{n-k}$ for every $s \geq 3$.
\end{theorem}

\begin{proof}
Since $\mathcal{C}$ has $q^\ell$ elements, $|\mathcal{F}| = q^\ell$. Each $U \in \mathcal{G}$ is an $\ell$-space, so each $U^\perp$ is a $(n-\ell) = k$-space.

A set-like $3$-sunflower-free family is automatically set-like $s$-sunflower-free for $s \geq 3$ (any $s$-sunflower contains a $3$-sub-sunflower with the same kernel). It therefore suffices to show that $\mathcal{F}$ is set-like $3$-sunflower-free.

By Proposition~\ref{prop:sum-complement}, for any $U, V, W \in \mathcal{F}$ corresponding to $A, B, C \in \mathcal{C}$,
\[
\dim(U \cap V) = n - \dim(U^\perp + V^\perp), \qquad \dim(U \cap V \cap W) = n - \dim(U^\perp + V^\perp + W^\perp).
\]

\textit{Pairwise intersection.} For distinct $A, B \in \mathcal{C}$, row reduction gives
\[
\mathrm{rank}\begin{pmatrix} I & A & [A^2]_1 & \mathbf{0} \\ I & B & [B^2]_1 & \mathbf{0} \end{pmatrix}
= \mathrm{rank}\begin{pmatrix} I & A & [A^2]_1 & \mathbf{0} \\ 0 & B - A & [B^2 - A^2]_1 & \mathbf{0} \end{pmatrix}
= 2\ell,
\]
since $B - A$ is invertible (Proposition~\ref{prop:matrix-rep}). Hence $\dim(U \cap V) = n - 2\ell$.

\textit{Triple intersection.} For pairwise distinct $A, B, C \in \mathcal{C}$,
\begin{align*}
\mathrm{rank}\begin{pmatrix} I & A & [A^2]_1 & \mathbf{0} \\ I & B & [B^2]_1 & \mathbf{0} \\ I & C & [C^2]_1 & \mathbf{0} \end{pmatrix}
&\;=\; \mathrm{rank}\begin{pmatrix} I & A & [A^2]_1 & \mathbf{0} \\ 0 & B - A & [B^2 - A^2]_1 & \mathbf{0} \\ 0 & C - A & [C^2 - A^2]_1 & \mathbf{0} \end{pmatrix} \\
&\;=\; \mathrm{rank}\begin{pmatrix} I & A & [A^2]_1 & \mathbf{0} \\ 0 & I & [B + A]_1 & \mathbf{0} \\ 0 & I & [C + A]_1 & \mathbf{0} \end{pmatrix} \\
&\;=\; \mathrm{rank}\begin{pmatrix} I & A & [A^2]_1 & \mathbf{0} \\ 0 & I & [B + A]_1 & \mathbf{0} \\ 0 & 0 & [C - B]_1 & \mathbf{0} \end{pmatrix} \;=\; 2\ell + 1.
\end{align*}
Here we used that $\mathcal{C}$ is a field (so closed under multiplication, inverses, and so that $(B-A)^{-1}(B^2 - A^2) = B + A$), together with the column-extraction identities $X[Y]_1 = [XY]_1$ and $[X]_1 + [Y]_1 = [X+Y]_1$ valid for all $X, Y \in \mathcal{C}$. Since $C \neq B$ implies $C - B$ is invertible, $[C-B]_1$ is a nonzero column and contributes rank $1$. Hence $\dim(U \cap V \cap W) = n - (2\ell + 1) = n - 2\ell - 1$.

\textit{Conclusion.} For pairwise distinct $U, V, W \in \mathcal{F}$,
\[
\dim(U \cap V \cap W) \;=\; n - 2\ell - 1 \;<\; n - 2\ell \;=\; \dim(U \cap V).
\]
If $\{U, V, W\}$ were a set-like $3$-sunflower with kernel $K$, then $K = U \cap V = U \cap W = V \cap W$, which forces $K = U \cap V \cap W$ and hence $\dim(U \cap V) = \dim(U \cap V \cap W)$, contradicting the strict inequality above. Therefore $\mathcal{F}$ contains no set-like $3$-sunflower.
\end{proof}

\end{document}
